% 04-graficas.tex — Sección 4: análisis gráfico (R-30). Las tres figuras
% obligatorias, dibujadas por pgfplots desde report/datos/*.dat.

\section{Análisis gráfico}
\label{sec:graficas}

Las tres figuras las dibuja LaTeX leyendo los archivos de
\texttt{report/datos/}, que genera \texttt{scripts/build\_report\_data.py} del
CSV y de \texttt{results/raw/}. Así cada punto es rastreable hasta su fila.

% p99-concurrencia.tex — Figura obligatoria 1 (R-30): P99 contra concurrencia.
%
% R-30.1 (decisión de Alix, al aprobar F7-e): la figura muestra LAS DOS VISTAS, no
% una nota al pie. Con c=50 en los perfiles B y C, hasta el 91 % de las consultas
% de la ventana son la primera consulta en frío de cada cliente (F7.3), así que
% parte de la subida de la curva es arranque en frío y no latencia del motor.
%   - línea continua: el P99 de la columna del CSV (lo que mide el manual);
%   - línea punteada: el P99 recalculado sin la primera consulta de cada cliente,
%     desde results/raw/, que es la vista en régimen.
% Donde el .dat trae `nan` no hay vista en régimen porque TODAS las consultas de
% la ventana eran primeras consultas: la línea punteada se interrumpe ahí
% (unbounded coords=jump, en el preámbulo).
%
% Los cuatro paneles van escritos uno a uno, SIN un macro con parámetros: dentro
% de un \newcommand, los '#' del cuerpo chocan con el analizador de claves de
% pgfplots ("Illegal parameter number in definition of \pgfkeys@..."). Es más
% largo, pero no depende de cómo pgfkeys reinterprete el cuerpo.

\begin{figure}[htbp]
  \centering
  \begin{tikzpicture}
    \begin{axis}[
        title={MongoDB · Q2 · perfil A}, title style={font=\small},
        xlabel={Concurrencia (clientes)}, ylabel={P99 (ms)},
        xmode=log, log basis x=10, xtick={1,10,50}, xticklabels={1,10,50},
        ymin=0, width=0.47\linewidth, height=5cm,
        legend pos=north west, legend cell align=left,
      ]
      \addplot[color=blue!60!black, mark=*, thick] table [x=concurrencia, y=p99_csv] {datos/p99_mongodb_Q2_A.dat};
      \addplot[color=blue!60!black, mark=o, densely dotted, thick]
        table [x=concurrencia, y=p99_sin_frias] {datos/p99_mongodb_Q2_A.dat};
      \legend{P99 del CSV, Sin la 1.\textsuperscript{a} en frío}
    \end{axis}
  \end{tikzpicture}\hfill
  \begin{tikzpicture}
    \begin{axis}[
        title={PostgreSQL · Q2 · perfil A}, title style={font=\small},
        xlabel={Concurrencia (clientes)}, ylabel={P99 (ms)},
        xmode=log, log basis x=10, xtick={1,10,50}, xticklabels={1,10,50},
        ymin=0, width=0.47\linewidth, height=5cm,
        legend pos=north west, legend cell align=left,
      ]
      \addplot[color=red!60!black, mark=*, thick] table [x=concurrencia, y=p99_csv] {datos/p99_postgresql_Q2_A.dat};
      \addplot[color=red!60!black, mark=o, densely dotted, thick]
        table [x=concurrencia, y=p99_sin_frias] {datos/p99_postgresql_Q2_A.dat};
    \end{axis}
  \end{tikzpicture}

  \begin{tikzpicture}
    \begin{axis}[
        title={MongoDB · Q2 · perfil C}, title style={font=\small},
        xlabel={Concurrencia (clientes)}, ylabel={P99 (ms)},
        xmode=log, log basis x=10, xtick={1,10,50}, xticklabels={1,10,50},
        ymin=0, width=0.47\linewidth, height=5cm,
        legend pos=north west, legend cell align=left,
      ]
      \addplot[color=blue!60!black, mark=*, thick] table [x=concurrencia, y=p99_csv] {datos/p99_mongodb_Q2_C.dat};
      \addplot[color=blue!60!black, mark=o, densely dotted, thick]
        table [x=concurrencia, y=p99_sin_frias] {datos/p99_mongodb_Q2_C.dat};
    \end{axis}
  \end{tikzpicture}\hfill
  \begin{tikzpicture}
    \begin{axis}[
        title={PostgreSQL · Q2 · perfil C}, title style={font=\small},
        xlabel={Concurrencia (clientes)}, ylabel={P99 (ms)},
        xmode=log, log basis x=10, xtick={1,10,50}, xticklabels={1,10,50},
        ymin=0, width=0.47\linewidth, height=5cm,
        legend pos=north west, legend cell align=left,
      ]
      \addplot[color=red!60!black, mark=*, thick] table [x=concurrencia, y=p99_csv] {datos/p99_postgresql_Q2_C.dat};
      \addplot[color=red!60!black, mark=o, densely dotted, thick]
        table [x=concurrencia, y=p99_sin_frias] {datos/p99_postgresql_Q2_C.dat};
    \end{axis}
  \end{tikzpicture}

  \caption[P99 contra concurrencia, en dos vistas]{%
    \textbf{P99 contra concurrencia, en dos vistas (R-30.1).} La línea continua es
    el P99 tal como lo mide el manual (columna del CSV). La punteada es el mismo
    P99 \textbf{sin la primera consulta de cada cliente}, recalculado desde
    \texttt{results/raw/}. La separación entre las dos líneas es el arranque en
    frío: el calentamiento de 60~s del manual no alcanza con \mbox{c=50} en los
    perfiles B y C, donde hasta el 91~\% de la ventana son primeras consultas
    (F7.3). En MongoDB · perfil C · \mbox{c=50} la línea punteada se interrumpe
    porque \emph{todas} las consultas de la ventana eran primeras consultas y no
    existe un P99 en régimen que dibujar. \textbf{Leer solo la línea continua
    sobreestima la latencia del motor en esos puntos.} Datos:
    \texttt{report/datos/p99\_*.dat}.}
  \label{fig:p99-concurrencia}
\end{figure}


La figura~\ref{fig:p99-concurrencia} lleva \textbf{dos series por panel} y no una
nota al pie, porque la diferencia entre ellas no es un detalle: el calentamiento
de 60~s que fija el manual no alcanza cuando 50 clientes empiezan a la vez con la
caché recién limpiada. Se midió una por una la primera consulta de cada cliente
y, con \mbox{c=50}, \textbf{no cabe en los 60~s en 14 de las 18 combinaciones};
solo cabe en el perfil~A con \texttt{Q2} (de 39 a 52~s). En los perfiles B y C
llega a 110~s, y en algunas celdas ninguna consulta terminó siquiera en los 110~s
que dura toda la corrida de prueba.

El caso extremo aparece en la figura: en \textbf{MongoDB, perfil~C,
\texttt{Q2}, \mbox{c=50}} las \num{150} consultas que terminaron dentro de la
ventana son \emph{todas} la primera de su cliente, así que no existe un P99 «en
régimen» que dibujar y la serie punteada se interrumpe. Decir que el P99 de esa
celda es de 228~s sin esta advertencia sería atribuirle al motor el costo de
arrancar en frío.

% throttling-qps.tex — Figura obligatoria 2 (R-30): throttling contra throughput.
%
% Lo que hay que saber para leer el eje (medido en F6.2, declarado en SPEC §4):
% `cpu_throttled_seconds` es la SUMA por CPU del tiempo que el cgroup tuvo hilos
% listos y no pudo correr. Por eso puede superar la duración de la ventana: con 4
% CPUs y una cuota de 1, cada segundo de reloj acumula hasta 4 de throttling. El
% eje NO se acota a 300 s.
%
% Además (F6.3): parte del throttling de MongoDB en el perfil C ocurre en reposo,
% sin benchmark (~27 s por ventana de 300 s). La línea de referencia lo marca.
%
% Un archivo de datos por motor y concurrencia (report/datos/throttling_*.dat),
% para no depender de filtros de pgfplots.

\begin{figure}[htbp]
  \centering
  \begin{tikzpicture}
    \begin{axis}[
        xlabel={Throttling de CPU en la ventana de 300~s (s)},
        ylabel={Throughput (consultas/s)},
        ymode=log, log basis y=10,
        xmin=0,   % el throttling no puede ser negativo
        enlarge x limits={upper, value=0.05},
        legend pos=south west, legend cell align=left, legend columns=2,
        height=7.5cm, width=\linewidth,
      ]
      \addplot[only marks, mark=*,       mark size=1.6pt, blue!65!black]
        table [x=throttling_s, y=qps] {datos/throttling_mongodb_c1.dat};
      \addplot[only marks, mark=*,       mark size=2.6pt, blue!65!black]
        table [x=throttling_s, y=qps] {datos/throttling_mongodb_c10.dat};
      \addplot[only marks, mark=*,       mark size=3.6pt, blue!65!black]
        table [x=throttling_s, y=qps] {datos/throttling_mongodb_c50.dat};
      \addplot[only marks, mark=square*, mark size=1.6pt, red!65!black]
        table [x=throttling_s, y=qps] {datos/throttling_postgresql_c1.dat};
      \addplot[only marks, mark=square*, mark size=2.6pt, red!65!black]
        table [x=throttling_s, y=qps] {datos/throttling_postgresql_c10.dat};
      \addplot[only marks, mark=square*, mark size=3.6pt, red!65!black]
        table [x=throttling_s, y=qps] {datos/throttling_postgresql_c50.dat};
      \legend{MongoDB c=1, MongoDB c=10, MongoDB c=50,
              PostgreSQL c=1, PostgreSQL c=10, PostgreSQL c=50}
      % F6.3: línea base de throttling en reposo de MongoDB en el perfil C.
      \draw[dashed, gray!70] ({axis cs:26.7,0}|-{rel axis cs:0,0})
                          -- ({axis cs:26.7,0}|-{rel axis cs:0,1});
      \node[gray!70, font=\tiny, anchor=south west, rotate=90]
        at ({axis cs:26.7,0}|-{rel axis cs:0,0.05}) {MongoDB · C en reposo (F6.3)};
    \end{axis}
  \end{tikzpicture}
  \caption[Throttling contra throughput]{%
    \textbf{Throttling de CPU contra throughput.} Cada punto es una celda (media de
    sus 3 repeticiones); el tamaño de la marca crece con la concurrencia.
    \textbf{El eje horizontal supera los 300~s de la ventana} porque
    \texttt{cpu\_throttled\_seconds} suma el tiempo de cada CPU en la que el cgroup
    tenía hilos listos (medido en F6.2, declarado en SPEC §4); el máximo observado
    es de \num{2086}~s. La línea vertical marca el throttling que MongoDB acumula
    \textbf{en reposo} en el perfil C, sin benchmark (F6.3): a su izquierda no hay
    contención atribuible a la carga. Datos:
    \texttt{report/datos/throttling\_*.dat}.}
  \label{fig:throttling-qps}
\end{figure}


% curva-rss.tex — Figura obligatoria 3 (R-30): curva de memoria durante la ventana.
%
% Dos series por panel, y la diferencia entre ellas es el hallazgo F6-A3:
%   anon        = lo que mide la columna `max_rss_mb` del CSV;
%   working_set = memory.current − inactive_file, que SÍ incluye la memoria
%                 compartida (`shared_buffers`) de PostgreSQL.
% Los `shared_buffers` son `shmem`, no `anon`, así que con `max_rss_mb` PostgreSQL
% sale artificialmente eficiente (SPEC §4 y A11). M3 se interpreta con working_set.
%
% Fuente: la serie de 1 s del cgroup que el runner guarda en results/raw/
% (serie_cgroup_1s), es decir las MISMAS lecturas de las que sale la columna.

\newcommand{\rssPanel}[4]{%
  % #1 archivo .dat · #2 título · #3 color · #4 "leyenda" para mostrarla
  \begin{tikzpicture}
    \begin{axis}[
        title={#2}, title style={font=\small},
        xlabel={Segundo de la ventana de medición}, ylabel={Memoria (MiB)},
        xmin=0, xmax=300, ymin=0,
        width=0.47\linewidth, height=5cm,
        legend pos=south east, legend cell align=left,
      ]
      \addplot[color=#3, thick, no marks] table [x=t_s, y=working_set_mib] {#1};
      \addplot[color=#3!55, thick, densely dashed, no marks]
        table [x=t_s, y=anon_mib] {#1};
      \addplot[color=gray!60, thin, dotted, no marks] table [x=t_s, y=shmem_mib] {#1};
      \ifthenelse{\equal{#4}{leyenda}}%
        {\legend{\texttt{working\_set}, \texttt{anon}, \texttt{shmem}}}{}%
    \end{axis}
  \end{tikzpicture}%
}

\begin{figure}[htbp]
  \centering
  \rssPanel{rss_mongodb_A_Q2_c50.dat}{MongoDB · perfil A · Q2 · c=50}{blue!60!black}{leyenda}\hfill
  \rssPanel{rss_postgresql_A_Q2_c50.dat}{PostgreSQL · perfil A · Q2 · c=50}{red!60!black}{}

  \rssPanel{rss_mongodb_C_Q2_c50.dat}{MongoDB · perfil C · Q2 · c=50}{blue!60!black}{}\hfill
  \rssPanel{rss_postgresql_C_Q2_c50.dat}{PostgreSQL · perfil C · Q2 · c=50}{red!60!black}{}

  \caption[Curva de memoria durante la ventana]{%
    \textbf{Memoria del contenedor durante la ventana, muestreada cada segundo.}
    La línea discontinua (\texttt{anon}) es exactamente lo que resume la columna
    \texttt{max\_rss\_mb}; la continua (\texttt{working\_set}) incluye además la
    memoria compartida. \textbf{En PostgreSQL la distancia entre las dos es
    \texttt{shared\_buffers}}: son \texttt{shmem}, no \texttt{anon}, así que
    \texttt{max\_rss\_mb} no cuenta su búfer principal y lo hace parecer
    artificialmente eficiente frente a MongoDB, cuyo caché de WiredTiger sí es
    \texttt{anon} (F6-A3, SPEC §4 y A11). Por eso M3 se interpreta con
    \texttt{working\_set}. Fuente: serie de 1 s del cgroup en
    \texttt{results/raw/}, las mismas lecturas de las que sale la columna.}
  \label{fig:curva-rss}
\end{figure}


La figura~\ref{fig:curva-rss} es la que obliga a tener cuidado con la columna
\texttt{max\_rss\_mb}, y por eso lleva dos series: véase la
sección~\ref{sec:memoria}.
